% Generated draft: requires mathematical and source review.
\documentclass[11pt]{article}
\usepackage[margin=25mm]{geometry}
\usepackage{fontspec}
\setmainfont{Latin Modern Roman}
\usepackage{amsmath,amssymb,mathtools,unicode-math}
\setmathfont{Latin Modern Math}
\usepackage{enumitem,needspace,xurl,hyperref}
\setlength{\parindent}{0pt}
\setlength{\parskip}{5pt}
\setlength{\emergencystretch}{3em}
\newcommand{\N}{\mathbb{N}}
\newcommand{\Z}{\mathbb{Z}}
\newcommand{\Q}{\mathbb{Q}}
\newcommand{\R}{\mathbb{R}}
\newcommand{\C}{\mathbb{C}}
\newcommand{\F}{\mathbb{F}}
\newcommand{\A}{\mathbb{A}}
\newcommand{\PP}{\mathbb{P}}
\newcommand{\E}{\mathbb{E}}
\newcommand{\eps}{\varepsilon}
\newcommand{\dd}{\,\mathrm{d}}
\begin{document}
\section*{UnsolvedMath ID: 30003034}
\textbf{Source problem code:} catalogue entry
\subsection{Definitions and mathematical statement}
Put \(R=\Z[1/2]=\{a/2^k:a\in\Z,\ k\in\Z_{\geq 0}\}\) and \(\F_2=\Z/2\Z\). For a positive integer \(n\), let \(GL_n(R)\) be the group of invertible \(n\times n\) matrices over \(R\), and interpret \(H^*(GL_n(R);\F_2)\) as the graded group cohomology algebra of this discrete group with trivial coefficients \(\F_2\).

Write \(\mathcal G_{n,2}(R)\) solely as a name for the object called ``Grothendieck's mod-2 Chern-class algebra'' in the supplied record. The record refers to a natural morphism
\[
\gamma_{n,R}:\mathcal G_{n,2}(R)\longrightarrow H^*(GL_n(R);\F_2).
\]
It does not give the construction of \(\mathcal G_{n,2}(R)\), its generators and relations, or the definition of \(\gamma_{n,R}\). Thus this notation preserves the source's intended objects without supplying an unverified construction.

The problem is to determine exactly which integers \(n\) in the intended range of matrix sizes make \(\gamma_{n,\Z[1/2]}\) an isomorphism of graded algebras. For positive matrix sizes, this asks to determine
\[
\left\{n\in\Z_{\geq 1}:\gamma_{n,\Z[1/2]}\text{ is an isomorphism of graded algebras}\right\}.
\]
The source says only ``integers \(n\)''; it does not specify whether \(n=0\) is included or provide a convention for that case. A fully self-contained formulation also requires the missing definitions of the source algebra and comparison morphism.

The supplied dated curation reports that the comparison is an isomorphism through \(n=3\) and fails for \(n\geq 14\), while the intermediate range is not settled by the cited results. These are supplied literature qualifications, not independently verified mathematical assertions here; they do not establish that every intermediate case remains open.
\subsection{Short English statement}
Determine precisely which matrix sizes make the source's natural comparison from Grothendieck's mod-2 Chern-class algebra to the mod-2 cohomology of \(GL_n(\Z[1/2])\) an isomorphism. The source algebra and map require explicit definitions, and inclusion of size zero is unspecified. The supplied classification is labelled partially solved and requires independent mathematical review.
\subsection{Sources}
\begin{itemize}
\item Supplied local UnsolvedMath registry record, permanent ID \(30003034\), title ``Chern Class Comparison for Arithmetic General Linear Groups,'' source problem code \texttt{catalogue entry}, in the Oberwolfach Reports Open Problems collection. The record gives proposed year 2016 and names no individual proposer.
\item \emph{Lattices and Applications in Number Theory} (2016), Oberwolfach Reports, Report 3/2016, pp.\ 87--154, DOI \texttt{10.4171/owr/2016/3}. Supplied report provenance: workshop reports and their open-problem discussions. \url{https://doi.org/10.4171/owr/2016/3}.
\item Dwyer, \emph{Homology isomorphisms between algebraic groups made discrete}. Supplied literature reference: \url{https://doi.org/10.1112/blms/25.2.145}.
\item Literature assessment dated \texttt{2026-08-21}: supplied dated curation, not independently verified. It labels the problem partially solved, reports isomorphism through \(n=3\) and failure for \(n\geq14\), and says that the intermediate range is not settled by the cited results. The record explicitly calls this dated literature triage derived from supplied OWR curation and requests independent verification of time-sensitive status and cited claims.
\item The supplied statement-status label is \texttt{corrected\_verified}. Its accompanying qualification states that only safe notation, typography, and TeX-delimiter cleanup was applied, with hypotheses, quantifiers, and subquestions preserved. This label is not independent verification of the mathematics in this draft.
\end{itemize}
\end{document}
